% =====================================================================
%  Section 5: Experiments and Numerical Results  --  REVISED for M-1
%  Grid: 9 alpha x 3 networks x 10 seeds x 2 topics = 540 runs
%  Supersedes the T=35 / 180-run configuration.
%  ALL CONTENT IN THIS FILE IS NEW -> rendered in green.
% =====================================================================
\begingroup\color{revgreen}

\subsection{Dataset Overview}
\label{sec:dataset}

We source data from the \texttt{r/politics} subreddit via PushShift's
\citep{baumgartner2020pushshift} April 2019 dump, containing 138.5 million comments.
After streaming decompression and filtering (non-deleted, non-removed, body length
$>$ 20 characters), 122.5 million active comments remain (deletion rate 11.56\%).
Semantic clustering (Section~\ref{sec:corpus}) identifies 1{,}184 topic clusters; we
focus on two politically contested topics: \emph{gun control} ($n = 3{,}581$) and
\emph{abortion} ($n = 1{,}618$).

\paragraph{Recalibration of the stance scale.}
Our earlier configuration scored comments with a RoBERTa regressor trained directly on
the IBM Argument Quality 30K corpus \citep{gretz2020argq}. Auditing that regressor
before the present experiments revealed a defect that invalidates any claim of a
\emph{continuous} stance scale: the supervising label \texttt{stance\_WA} takes only two
distinct values in the training split, so the fitted head behaved as a saturating
classifier rather than a regressor. Empirically, 69.7\% of scored comments fell in
$|\hat{y}| > 0.9$, and the mapping from human-judged stance strength to $\hat{y}$ was
non-monotone in the interior of the range. A regressor that places seven of ten comments
at the endpoints cannot express the gradations the model requires.

We therefore rebuilt the supervision using Best--Worst Scaling to induce a graded target
from pairwise comparisons, and re-fine-tuned the regression head against it. After
recalibration, group means are strictly monotone in human-judged stance strength, the
correlation with the graded target is $+0.926$, and the saturation rate falls from
69.7\% to 6.5\% (neutral-group mean $-0.020$). All initial beliefs and all simulations
reported here use the recalibrated scorer; results obtained under the saturating scorer
are not comparable to them and are not reported.

Figure~\ref{fig:stance-dist} shows the recalibrated stance distributions. Gun control
exhibits a near-neutral mean with high variance, consistent with a bimodal and divided
public; abortion displays a more concentrated opposing consensus. The resulting
initial-stance split among the $N = 50$ sampled agents is 54/46 for gun control and
68/32 for abortion --- an asymmetry that Section~\ref{sec:results} shows to be
consequential.

\paragraph{The two topics are not independent replicates.}
The 50 sampled agents share \emph{identical} personas (name, age, gender, education,
five personality traits) and \emph{identical} stubbornness coefficients $\rho_i$ across
the two topics; only the initial stance differs. Measured across the population,
$\mathrm{corr}(s^{\text{gun}}, s^{\text{abortion}}) = +0.881$. This is a deliberate
control: it allows any topic effect to be attributed to the stance distribution rather
than to a different population of agents. It also means the two topics must be compared
with \emph{paired} tests and must not be pooled as $2\times$ independent samples
(Section~\ref{sec:stats}). We note as a modeling simplification --- not an empirical
finding --- that the same individual is assigned the same stubbornness on both issues.

\subsection{Simulation Configuration}
\label{sec:config}

All simulations use $N = 50$ agents and $K = 5$ out-neighbors per agent. The mixing
parameter is varied across
$\alpha \in \{0, 0.125, 0.25, 0.375, 0.5, 0.625, 0.75, 0.875, 1\}$, spanning the full
range from a purely LLM-driven population ($\alpha = 0$) to a purely cost-minimizing
population ($\alpha = 1$).

\paragraph{Initial topology as an experimental factor.}
Rather than fixing a single scale-free initialization, we treat the initial topology as
a manipulated factor with three levels: Barabasi--Albert (BA, $m=2$)
\citep{barabasi1999emergence}, Erdos--Renyi (ER), and Watts--Strogatz (WS). This
directly tests whether echo chambers and efficiency loss depend on pre-existing
structure or are produced by the K-NN rewiring rule itself. \emph{The three families are
generated with matched mean degree} --- 3.84 (BA), 3.84 (ER), 4.00 (WS), a maximum
deviation of 4.17\% --- so that any observed difference is attributable to structure
rather than to edge count. Their clustering coefficients remain clearly distinct
(0.179 / 0.079 / 0.390), preserving the contrast of interest. Ten independent
topologies (seeds 1--10) are generated per family.

\paragraph{Grid.}
The full design is
$9\ (\alpha) \times 3\ (\text{network}) \times 10\ (\text{seed}) \times 2\ (\text{topic})
= \mathbf{540}$ runs, with $n = 60$ runs at every $\alpha$ and $n = 30$ at every
$(\alpha, \text{topic})$ cell. The completed grid contains no missing cells.

\paragraph{Simulation horizon: dynamic, not fixed.}
Our earlier configuration ran every simulation for a fixed $T = 35$ steps. That horizon
is not defensible under the convergence criterion developed in
Section~\ref{sec:convergence}: measured mean settling time at $\alpha = 0$ is 52.2 steps,
with individual runs requiring up to 115. A fixed $T = 35$ truncates the majority of
low-$\alpha$ runs before they reach an attractor, so metrics read at $t = 35$ would
describe a transient rather than a converged state. We therefore replace the fixed
horizon with the dynamic stopping rule of Section~\ref{sec:stopping}: each run terminates
\texttt{post\_window} $= 10$ steps after its own detected convergence time
$t_{\mathrm{conv}}$, subject to a hard cap $T_{\max} = 120$. Across the grid, 538 of 540
runs (99.63\%) converged within the cap; the two exceptions are recorded as
\texttt{attractor} $=$ \texttt{none} and retained rather than discarded.

\paragraph{Belief representation.}
Type-L agents no longer emit a numeric belief on a five-point scale. They produce text
only, and their position on $[-1,+1]$ is obtained by scoring that text with the same
recalibrated RoBERTa regressor used for the Reddit corpus (Section~\ref{sec:dataset}).
Agent positions and the empirical distribution are therefore measured on one common
ruler. This change is not cosmetic: under the five-point scale, quantization forced ties
at the K-NN cutoff for 50 of 50 agents, so neighbor selection was determined almost
entirely by tie-breaking rather than by opinion. Ties are now broken by seeded random
choice, reducing the correlation between in-degree and agent index from $-0.22$ to
$-0.037$.

\paragraph{Determinism.}
Deterministic decoding (temperature $= 0$) is used throughout and the vLLM server is
initialized with a fixed seed. The numerical ($\alpha = 1$) path was verified to be
bit-for-bit reproducible across repeated runs.

\subsection{Evaluation Metrics}
\label{sec:metrics}

Unless stated otherwise, cross-run comparisons are read at
$t_{\mathrm{conv}} + \texttt{post\_window}$ --- each run's own converged state --- rather
than at a common step index, since runs now differ in length. We additionally record
values at the fixed indices $t = 20$ and $t = 35$ for reference.

\paragraph{Polarization ($P_z$).}
Following \citet{chitra2020filter}, the variance of expressed opinions,
$P_z^{(t)} = \frac{1}{N}\sum_i \bigl(z_i^{(t)} - \bar{z}^{(t)}\bigr)^2$.

\paragraph{Modularity: $Q$ and the null-model correction $Q_{\mathrm{norm}}$.}
Louvain community detection \citep{blondel2008fast} on the undirected projection of
$\mathcal{G}^{(t)}$ yields a modularity score $Q^{(t)}$ \citep{newman2004finding}.
Bare $Q$ has no absolute meaning: \citet{guimera2004modularity} showed that random
graphs exhibit non-zero modularity from fluctuations alone, and our K-NN rule imposes a
constant out-degree $K$, a structural constraint that itself induces modularity. We
therefore report $Q$ against a degree-preserving null model:
\[
Q_{\mathrm{norm}}(t) = Q_{\mathrm{obs}}(t) - \bar{Q}_{\mathrm{rand}}(t),
\qquad
z_Q(t) = \frac{Q_{\mathrm{obs}}(t) - \bar{Q}_{\mathrm{rand}}(t)}{\sigma_{\mathrm{rand}}(t)},
\]
where $\bar{Q}_{\mathrm{rand}}$ averages $R = 20$ randomizations of the undirected
projection by \texttt{double\_edge\_swap} (\texttt{nswap} $= 10|E|$,
\texttt{max\_tries} $= 1000|E|$), preserving the degree sequence. The null is recomputed
at every step.

We record a methodological caution. We first implemented the null with directed
three-cycle swaps; that randomization \emph{does not mix} on this graph family --- edge
overlap plateaued at 0.35 even at \texttt{nswap} $= 50|E|$, and $\bar{Q}_{\mathrm{rand}}$
jumped from 0.239 to 0.654 between adjacent steps whose degree sequences were nearly
identical, while within-step replicate variability remained at $\mathrm{sd} \le 0.015$
(spurious precision). A null model must not depend that sharply on its input. The
undirected \texttt{double\_edge\_swap} reduces edge overlap to 0.13 and yields
$\bar{Q}_{\mathrm{rand}}$ stable to $\pm 0.005$ over the same steps. Swap failures are
counted and reported; a failed randomization is never silently replaced by a
configuration-model fallback. Across the grid, \texttt{swap\_fail} $= 0$.

We report bare $Q$ alongside $Q_{\mathrm{norm}}$ for comparability with prior work, but
all substantive claims use $Q_{\mathrm{norm}}$.

One limitation the null model cannot address must be stated explicitly: it controls for
the degree sequence but not for the \emph{resolution limit}. The bound of
\citet{fortunato2007resolution} is approximately $\sqrt{2m} \approx 22$ at $K = 5$,
$N = 50$, whereas the smallest communities we observe contain 5--7 nodes. Louvain may
therefore merge communities that should remain separate, and $Q_{\mathrm{norm}}$ does
not correct for this.

\paragraph{Price of Anarchy and its decomposition.}
PoA is defined as in Section~\ref{sec:poa}, with both $C(z^{(t)})$ and $C(z^*_t)$
recomputed on the same post-rewiring graph. Because PoA is an aggregate, it cannot
distinguish a population that is inefficient because it \emph{disagrees persistently}
from one that is inefficient because it \emph{abandons its own position}. The social
cost already separates into exactly these two terms:
\[
C(z) = \underbrace{\sum_i \sum_{j \in \mathcal{N}(i)} (z_i - z_j)^2}_{\text{disagreement}}
     \; + \; \underbrace{\sum_i \rho_i K (z_i - s_i)^2}_{\text{conformity}} .
\]
We report both components. The decomposition is recomputed \emph{post hoc} from stored
per-step agent states and edge sets, leaving the metric schema untouched; every run's
recomputed PoA is cross-validated against the value logged during simulation (median
relative error $< 2\%$), and any run failing validation is reported rather than silently
included. All 540 runs passed.

\paragraph{Structural convergence metrics.}
$C^{\mathrm{out}}(t)$, $dS(t)$, $d_L(t)$, DeltaCon$(t)$ and partition ARI are defined in
Section~\ref{sec:convergence} and written at every step.

\subsection{Statistical Analysis}
\label{sec:stats}

Variances are markedly unequal across $\alpha$ --- the standard deviation of PoA falls by
two orders of magnitude between $\alpha = 0$ and $\alpha = 1$ --- so comparisons between
$\alpha$ levels use \textbf{Welch's $t$-test} with Benjamini--Hochberg FDR control over
the family of adjacent-$\alpha$ comparisons; we report $q$-values.

Claims of \emph{no difference} are not supported by $p > 0.05$. Where we assert
equivalence --- specifically for the comparison among initial topologies --- we use
\textbf{two one-sided tests (TOST)} with the equivalence margin fixed in advance at
$\delta = 0.2$ PoA units. We flag a scale limitation of this choice: $\delta = 0.2$ is a
17.5\% relative margin at $\alpha = 1$ (PoA $\approx 1.14$) but only 3.6\% at
$\alpha = 0$ (PoA $\approx 5.56$), so the same absolute margin is five times stricter at
the low-$\alpha$ end. We report the consequence of this in Section~\ref{sec:results}
rather than selecting a margin after seeing the data.

Comparisons \emph{between topics} use paired tests, for the reason given in
Section~\ref{sec:dataset}. All confidence intervals are 95\%.

\subsection{Implementation Details}
\label{sec:impl}

Type-L agents are powered by Phi-4 (14.7B parameters, AWQ 4-bit quantization)
\citep{abdin2024phi4}, served locally via vLLM \citep{kwon2023vllm} on a single
RTX 4090. Type-C agents are closed-form FJ updates requiring no neural inference.
Network operations use NetworkX and \texttt{python-louvain}. The full simulation
codebase is available at
\url{https://github.com/Evan-Jiamg/Echo-Chamber-Simulation}.

\paragraph{One call per agent per step.}
Each Type-L agent makes exactly one LLM call per step, emitting a structured record with
fields generated in the order \texttt{(reasoning, opinion, memory)}. The ordering is
load-bearing: fields are generated in declaration order, so declaring \texttt{reasoning}
first is what makes the chain of thought precede the stated position rather than
rationalize it after the fact. The prompt requires the agent to name the specific
arguments it heard this round, judge which are strong, relate them to its original
stance, and only then state the resulting belief.

\paragraph{Memory.}
Memory is maintained by the model itself: at each step the agent rewrites a
$\le$120-word memory, instructed to carry forward what remains relevant, incorporate what
is new, and drop what is stale. This replaces a hand-written compression routine that we
found to be defective --- it retained the \emph{first} $N$ words rather than the most
recent, with accumulated history placed at the front of the concatenated string, so
historical memory \emph{froze after step 1} and no later information could enter. We
verified on the completed grid, not merely by code review, that no agent's memory is
frozen: 0 of 50 agents in every run. Neighbor opinions are now passed through at up to
60 words each, replacing a 15-word truncation that discarded roughly 70\% of a typical
40--60-word opinion.

\paragraph{Throughput and stability.}
Guided (schema-constrained) decoding is disabled. Its finite-state mask serializes on the
CPU and costs a factor of ten in throughput, while measured parse success without it is
200/200; leaving it enabled would have extended the grid from days to more than a month.
Parse failures are nonetheless counted per run and were zero throughout. Prefix caching
contributes a further 27\% (hit rate 54--82\%). The complete grid required approximately
40 GPU-hours. The vLLM engine raised sporadic CUDA illegal-memory-access faults (5
occurrences over the grid); a supervisor process restarted the server and resumed from
the last completed run, with 40--68 s of downtime per event and no data loss. Because a
parse failure would otherwise be indistinguishable from an agent deliberately holding its
position, a run aborts if two consecutive steps fail to parse entirely.

\paragraph{Data integrity.}
Every run writes six artifacts: per-step metrics (18 columns), a convergence summary,
full per-agent belief/opinion/reasoning/memory traces, per-step edge sets, a metrics
JSON-lines log, and the Type-C/Type-L assignment. An audit over all 540 runs confirmed:
no missing or zero-byte files; 18-column schema conformance 540/540; CSV row count equal
to \texttt{steps\_run} in every run (no truncation); reasoning traces retained in full
(mean 591--848 characters per step); zero frozen memories; zero swap failures; and no
missing grid cells. The complete output totals 2.34 GB.

\endgroup
